\documentclass[10pt]{article}
\usepackage[margin=0.78in]{geometry}
\usepackage{amsmath,amssymb,amsthm,mathtools,microtype}
\usepackage[T1]{fontenc}
\usepackage{lmodern}
\usepackage{graphicx}
\usepackage[hidelinks]{hyperref}
\newtheorem{theorem}{Theorem}
\newtheorem{proposition}{Proposition}
\newcommand{\one}{\mathbf 1}
\newcommand{\R}{\mathbb R}
\newcommand{\cF}{\mathcal F}
\newcommand{\diag}{\operatorname{diag}}
\setlength{\parindent}{0pt}
\setlength{\parskip}{3.5pt}
\title{A family of counterexamples and an exact\\
signless-Laplacian resolvent formula}
\author{Automated mathematical discovery run \texttt{fa\_banach\_001}, lane 11}
\date{17 August 2026}

\begin{document}
\maketitle

\textbf{Status.} Likely valid counterexample to Conjecture 8.1 as printed,
and likely valid full answer to Problem 8.2 of Hillar--Lin--Wibisono
\cite{HLW}.  The stronger motivating version of the conjecture is not
settled.

\section*{Source and exact questions}

C. Hillar, S. Lin, and A. Wibisono, ``Inverses of symmetric, diagonally
dominant positive matrices and applications,'' arXiv:1203.6812, Section 8
\cite{HLW}.  The paper defines $A\ge B$ to mean that $A-B$ is entrywise
nonnegative, reserving $A\succeq B$ for Loewner order.  With
$S(\alpha,m)=\alpha I+m\one\one^{\mathsf T}$, Conjecture 8.1 says that every
symmetric diagonally dominant $0<J\le S(\alpha,m)$ obeys
\[
 \|J^{-1}\|_\infty\ge \|S(\alpha,m)^{-1}\|_\infty
 =\frac{\alpha+2m(n-1)}{\alpha(\alpha+mn)},                 \tag{1}
\]
with equality only for $J=S$.  Problem 8.2 asks for exact combinatorial
formulas for all coefficients of the rational functions
$(S+tP)^{-1}_{ij}$ when $P$ is a graph signless Laplacian.

\begin{center}
\includegraphics[width=.96\textwidth]{figures/open_problem_crop.png}
\end{center}

\clearpage
\section*{Proof intuition}

The printed conjecture has forgotten information present in its motivating
paragraph.  Entrywise domination allows one to keep every diagonal equal to
the diagonal of $S$ while making all off-diagonal entries tiny.  The resulting
matrix is nearly scalar, so its inverse norm is much smaller than that of
$S$.  This produces counterexamples for every parameter choice.

For the graph problem, factor
\[
 \alpha I+\ell\one\one^{\mathsf T}+tBB^{\mathsf T}=XX^{\mathsf T}
\]
using weighted coordinate, all-ones, and unsigned-incidence columns.
Cauchy--Binet expands the determinant.  Applying the same identity to the
complementary minors in the cofactor formula expands every inverse-entry
numerator.  Counting chosen edge columns reads off every power of $t$.

\section*{A universal counterexample family}

\begin{theorem}\label{thm:counterexample}
For every $n\ge3$ and every $\alpha,m>0$, Conjecture 8.1 as printed is false.
More precisely, for all sufficiently small $\varepsilon>0$,
\[
 J_\varepsilon=(\alpha+m-\varepsilon)I
                  +\varepsilon\one\one^{\mathsf T}        \tag{2}
\]
is positive entrywise, symmetric, diagonally dominant, and entrywise at most
$S(\alpha,m)$, but
$\|J_\varepsilon^{-1}\|_\infty<\|S(\alpha,m)^{-1}\|_\infty$.
\end{theorem}

\begin{proof}
The diagonal of $J_\varepsilon$ is $\alpha+m$, equal to that of $S$, and
its off-diagonal entries are $\varepsilon\le m$.  It is diagonally dominant
provided
\[
                  \alpha+m\ge(n-1)\varepsilon.             \tag{3}
\]
Thus every sufficiently small positive $\varepsilon$ satisfies all printed
hypotheses.  The inverse of $aI+b\one\one^{\mathsf T}$ is
\[
 \frac1aI-\frac{b}{a(a+nb)}\one\one^{\mathsf T}.
\]
Its diagonal inverse entries are positive and its off-diagonal entries are
negative.  Taking $a=\alpha+m-\varepsilon$ and $b=\varepsilon$ therefore
gives
\[
 \|J_\varepsilon^{-1}\|_\infty
 =\frac{\alpha+m+(2n-3)\varepsilon}
 {(\alpha+m-\varepsilon)(\alpha+m+(n-1)\varepsilon)}.      \tag{4}
\]
The right side tends to $1/(\alpha+m)$ as $\varepsilon\downarrow0$.  On the
other hand, the proposed lower-bound value in (1) satisfies
\[
 \|S^{-1}\|_\infty-\frac1{\alpha+m}
 =\frac{m(n-1)(\alpha+2m)}
 {\alpha(\alpha+mn)(\alpha+m)}>0.                          \tag{5}
\]
Continuity of (4), together with (3), completes the proof.
\end{proof}

For the smallest displayed instance, take $n=3$ and $\alpha=m=1$:
\[
 S=I+\one\one^{\mathsf T},\qquad
 J=\tfrac32I+\tfrac12\one\one^{\mathsf T}.
\]
Then $J\le S$ entrywise and is diagonally dominant, while
\[
 S^{-1}=I-\tfrac14\one\one^{\mathsf T},\quad
 J^{-1}=\tfrac23I-\tfrac19\one\one^{\mathsf T},\quad
 \|S^{-1}\|_\infty=\tfrac54>\tfrac79=\|J^{-1}\|_\infty.
\]

\section*{Why this does not settle the natural repair}

Immediately before Conjecture 8.1, the source assumes separate upper bounds
$J_{ij}\le m$ and $\Delta_i(J)\le\delta$ and observes that they imply
$J\le S(\alpha,m)$ for $\alpha=m(n-2)+\delta$.  The converse is false.  For
(2),
\[
 \Delta_i(J_\varepsilon)-\delta=(n-1)(m-\varepsilon)>0
 \qquad(0<\varepsilon<m).                                 \tag{6}
\]
The example therefore exploits precisely the missing dominance cap.  The
version retaining both motivating bounds remains open.

\section*{An exact answer to the graph resolvent problem}

Let $G=(V,E)$ be a loopless graph on $V=\{1,\ldots,n\}$, let $B$ be its
unsigned incidence matrix, and let $P=BB^{\mathsf T}$.  We write
\[
                   M(t)=\alpha I+\ell\one\one^{\mathsf T}+tP,
                   \qquad \alpha,\ell,t>0.                 \tag{7}
\]
Form the ground set $\Omega=V\sqcup\{\star\}\sqcup E$.  Associate columns
and weights
\[
\begin{array}{c|c|c}
q& a_q&w_q\\ \hline
v\in V&e_v&\alpha\\
\star&\one&\ell\\
e=\{u,v\}\in E&e_u+e_v&t.
\end{array}                                                \tag{8}
\]
Fix an order on $\Omega$.  For $F\subseteq\Omega$, let $A_F$ contain its
columns in that order and let $w(F)=\prod_{q\in F}w_q$.  For $|F|=n-1$ set
\[
                  \Delta_i(F)=\det[\,e_i\; A_F\,],         \tag{9}
\]
where $e_i$ is the first column.

\begin{theorem}\label{thm:resolvent}
Every entry of the resolvent (7) has the exact combinatorial expression
\[
 (M(t)^{-1})_{ij}=\frac{N_{ij}(t)}{D(t)},                  \tag{10}
\]
where
\begin{align}
 D(t)&=\sum_{\substack{F\subseteq\Omega\\|F|=n}}
              w(F)\det(A_F)^2,                            \tag{11}\\
 N_{ij}(t)&=\sum_{\substack{F\subseteq\Omega\\|F|=n-1}}
              w(F)\Delta_i(F)\Delta_j(F).                 \tag{12}
\end{align}
In particular, the coefficient of $t^k$ in either polynomial is obtained by
restricting its sum to subsets containing exactly $k$ members of $E$.
\end{theorem}

\begin{proof}
Let $X$ have the columns $\sqrt{w_q}a_q$ for $q\in\Omega$.  Equations (7)
and (8) give $M(t)=XX^{\mathsf T}$.  Cauchy--Binet applied to this product is
exactly (11).

For the numerator, the cofactor formula and Cauchy--Binet on the
$(n-1)\times(n-1)$ complementary minors give
\[
 \det M(t)\,(M(t)^{-1})_{ij}
 =(-1)^{i+j}\!\sum_{|F|=n-1}w(F)
   \det(A_F^{\widehat j})\det(A_F^{\widehat i}),           \tag{13}
\]
where a hat means that row is deleted.  Expansion of (9) along its first
column yields
$\Delta_i(F)=(-1)^{i+1}\det(A_F^{\widehat i})$.
Substitution in (13) proves (12).  Finally, only graph-edge weights contain
$t$, so grouping terms by $|F\cap E|$ gives the coefficient assertion.
\end{proof}

\section*{Rooted-forest interpretation}

The determinants in (11)--(12) are already finite combinatorial incidence
minors.  They also have a standard graph description.  If $\star\notin F$,
put $R=F\cap V$ and $H=(V,F\cap E)$.  Then $\det(A_F)\ne0$ exactly when
every component is one of:
\begin{itemize}
\item a tree containing exactly one root from $R$, contributing magnitude
      $1$; or
\item an odd-unicyclic component containing no root, contributing magnitude
      $2$.
\end{itemize}
Consequently $\det(A_F)^2=4^u$, where $u$ is the number of odd-unicyclic
components.  For a nonzero term with $\star\in F$, exactly one component is
instead an unrooted tree $T_0$; if its bipartition has sizes $p$ and $q$, its
additional squared factor is $(p-q)^2$.  The other components have the two
types above.

For completeness, this follows by splitting the incidence matrix into
connected blocks.  A rooted tree block has determinant magnitude $1$; an
unsigned unicyclic incidence block is nonsingular precisely for an odd cycle,
when its determinant has magnitude $2$.  For the unique unrooted tree, the
signed bipartition vector spans the left nullspace of its incidence matrix,
so adjoining $\one$ gives determinant magnitude $|p-q|$.  The numerator
minors have the same description after adjoining an implicit root at $i$;
the product in (12) retains the signs required off the diagonal.  Under the
source's self-loop convention, the formula itself extends verbatim by using
the incidence column $\sqrt2e_v$.

\section*{Verification, scope, and novelty}

The accompanying script exhausts all 74 labelled simple graphs on two,
three, or four vertices, at three positive $t$ values per graph.  It checks
(11)--(12) against direct determinants and inverses in all 222 cases, as well
as the exact $3\times3$ counterexample.  This is a sign/index audit, not part
of the proof.

A bounded search through 17 August 2026 used the source id/title, citation
trails, and close variants of ``Conjecture 8.1,'' ``signless Laplacian inverse
entries,'' and ``exact combinatorial formula $(S+tP)^{-1}$.''  It found
TU-subgraph formulas for signless-Laplacian characteristic polynomials and
specialized inverse formulas, but no answer to Conjecture 8.1 or Problem 8.2.
Because Theorem \ref{thm:resolvent} is an elementary all-minors
Cauchy--Binet identity, its novelty may be folklore; the claim is only
``apparently new within the bounded search.''

The source's final determinant conjecture is not new: Minghua Lin proved it
in arXiv:1212.1934 \cite{Lin}, and the proof appears in the authors' later
Hadamard-type paper \cite{HW}.  No originality is claimed for that status
identification.

\section*{Human-review recommendation}

First verify the source's entrywise convention and distinguish the printed
conjecture from its stronger motivating repair.  Then audit the deleted-row
order in (13); once the cofactor signs match (9), the resolvent formula is a
direct Cauchy--Binet identity.

\begin{thebibliography}{9}
\bibitem{HLW}
C. Hillar, S. Lin, and A. Wibisono,
``Inverses of symmetric, diagonally dominant positive matrices and
applications,'' arXiv:1203.6812 (2012).

\bibitem{Lin}
M. Lin, ``An upper bound for the determinant of a diagonally balanced
symmetric matrix,'' arXiv:1212.1934 (2012).

\bibitem{HW}
C. Hillar and A. Wibisono, ``A Hadamard-type lower bound for symmetric
diagonally dominant positive matrices,'' Linear Algebra Appl. 472 (2015),
135--141.
\end{thebibliography}

\end{document}
